% Resultados literales de la ejecucion del 13/08/2026.
% Fuente: docs/adr/adr-0004-base-vectorial.md, bloque escrito por el propio
% script entre los marcadores de RESULTADOS AUTOMATICOS.
% NO editar estas cifras a mano: si cambian, se vuelve a ejecutar
%   py scripts/experimento_vectordb.py
% y se copian de ahi. Corpus de 1.334 fragmentos de 384 dimensiones (1,95 MiB)
% y 50 preguntas. K = 10, 5 repeticiones por pregunta para la latencia.
\begin{table}[htbp]
  \centering
  \small
  \caption[Comparativa de bases de datos vectoriales]{\textcolor{borrador2}{Comparativa
      de bases de datos vectoriales sobre la colección completa y las 50 preguntas del
      conjunto de evaluación. Fuente: elaboración propia.}}
  \label{tab:comparativa_vectordb}
  \begin{tabular}{@{}lrrrrr@{}}
    \toprule
    \textbf{Candidata}   & \textbf{Fidelidad} & \textbf{Latencia} & \textbf{p90}
                         & \textbf{Construcción}                  & \textbf{Memoria}   \\
                         & \textbf{(U1)}      & \textbf{(U3)}     &
                         &                                        & \textbf{(U5)}      \\
    \midrule
    NumPy (línea base)   & \textbf{1,0000}    & \textbf{0,20}     & \textbf{0,24}
                         & \textbf{0,00}                          & \textbf{1,95}      \\
    ChromaDB 1.5.9       & 0,9980             & 1,43              & 1,68
                         & 1,81                                   & 26,12              \\
    LanceDB 0.37.1$^{\dagger}$ & \textbf{1,0000} & 7,35            & 8,41
                         & 0,30                                   & 22,71              \\
    Qdrant 1.19.0        & \textbf{1,0000}    & 5,62              & 16,61
                         & 1,51                                   & 149,10             \\
    \bottomrule
  \end{tabular}
  \begin{minipage}{0.95\textwidth}
    \vspace{4pt}
    \tiny
    \textbf{Notas:} las latencias y la construcción están en milisegundos y segundos
    respectivamente, y la memoria en MiB. \textbf{Fidelidad} es la proporción de vecinos que
    cada candidata devuelve y que coinciden con los de la búsqueda exacta, de modo que mide si
    el índice pierde resultados, no la calidad de la recuperación, que depende del modelo de
    incrustaciones y del troceado. \textbf{p90} es el percentil 90 de la latencia. NumPy no es
    una candidata: es la búsqueda exacta por fuerza bruta, que sirve de referencia de fidelidad
    y de línea base de velocidad. Las cuatro condiciones reciben los mismos vectores,
    incrustados una sola vez. En negrita, el mejor valor de cada columna.
    \textcolor{borrador2}{$^{\dagger}$~Candidata adoptada; el porqué, en el
      Apartado~\ref{secc:eleccion_lancedb}. Nótese que la elección no recae sobre la más rápida:
      la diferencia de latencia entre las dos que superan todos los umbrales queda dentro de la
      banda de indiferencia fijada antes de medir, y decide el desempate por simplicidad.}
  \end{minipage}
\end{table}
